\clearpage\fancyhead[L]{\small PALLAS.}\hypertarget{native-327-lower}{}\label{native-327-lower}
\newcommand{\ObservedDeclinationMark}{\hbox{\vrule width0pt height6.625000bp depth0pt\special{pdf:bcontent}\special{pdf:code q 0 g 0.333333 6.625000 m 0.333333 3.791667 l 0.166667 3.750000 l 0.125000 3.541667 l 0.000000 3.541667 l 0.000000 3.041667 l 0.125000 3.041667 l 0.166667 2.875000 l 0.333333 2.833333 l 0.333333 0.916667 l 0.166667 0.875000 l 0.166667 0.166667 l 0.291667 0.166667 l 0.333333 0.000000 l 0.833333 0.000000 l 0.875000 0.166667 l 1.041667 0.166667 l 1.041667 0.875000 l 0.833333 0.916667 l 0.833333 2.833333 l 1.041667 2.875000 l 1.083333 3.041667 l 3.000000 3.041667 l 3.041667 2.875000 l 3.541667 2.875000 l 3.583333 3.041667 l 3.916667 3.041667 l 3.916667 3.208333 l 4.083333 3.250000 l 4.083333 3.375000 l 4.250000 3.416667 l 4.250000 3.541667 l 4.083333 3.583333 l 4.083333 3.750000 l 1.416667 3.750000 l 1.375000 3.541667 l 1.083333 3.541667 l 1.041667 3.750000 l 0.833333 3.791667 l 0.833333 6.083333 l 1.041667 6.125000 l 1.041667 6.458333 l 0.875000 6.458333 l 0.833333 6.625000 l h f* Q}\special{pdf:econtent}\kern4.250000bp}}

\begin{center}\textbf{PALLAS.}\par\vspace{4pt}\hrule\vspace{4pt}\fontsize{10}{12}\selectfont Göttingische gelehrte Anzeigen. Stück 8 u. 9. Seite 73\,..\,77. 1811 Jan. 14.\par\vspace{4pt}\hrule\vspace{4pt}\end{center}

\par\indent Wir haben im 198. St. unserer Gel. Anz. vom vor. Jahre [S. 61 d. B.] eine Nachricht von einer Untersuchung des Hrn Prof. \textsc{Gauss} über die bisher beobachteten sechs Oppositionen der Pallas gegeben, welche sich nicht mehr durch eine rein elliptische Bewegung darstellen lassen. In der Ueberzeugung, dass die nun immer beträchtlicher werdenden Unterschiede eine Folge der Störungen sind, welche die Pallas besonders vom Jupiter erleidet, unternahm Hr. Prof. \textsc{Gauss} bereits im October des vorigen Jahrs die Berechnung der vom Jupiter herrührenden Störungen während der Jahre 1803..1811, theils um zu sehen, um wie viel sich nach Anbringung dieser Perturbationen die Oppositionen besser vereinigen lassen würden, als vorher, theils um dadurch die künftige Entwickelung der allgemeinen Theorie dieser Störungen vorzubereiten. Nach Vollendung dieser delicaten und schwierigen Arbeit fand sich ein so befriedigendes Resultat, als man kaum hätte erwarten dürfen: alle sechs Oppositionen liessen sich sowohl nach der Länge, als nach der Breite bis auf wenige Secunden genau vereinigen. Da nun gegen Ende des Jahrs 1810 Pallas wieder sichtbar werden musste, so schien es interessant, im voraus zu bestimmen, \SourceLeaf{328}wie viel die nach rein elliptischen Elementen berechnete Ephemeride für diese nächste Erscheinung durch die mit Rücksicht auf die Jupitersstörungen verbesserte Theorie des Planeten abgeändert werden würde, um so mehr, da der Einfluss des Jupiters bei der heliocentrischen Conjunction 1809 vorzüglich bedeutend gewesen war, und die Elemente sehr ansehnlich abgeändert hatte: vorzüglich war nemlich die Excentricität und mittlere Bewegung stark verkleinert, die mittlere Entfernung vergrössert, und die Apsidenlinie zurückgegangen. Es ergab sich hieraus, nach angestellter Rechnung, eine Abänderung der Ephemeride, weit grösser, als in irgend einem der vorhergehenden Jahre Statt gefunden hatte; die gerade Aufsteigung fiel gegen Ende des Octobers um \emph{drei Viertelgrad} kleiner aus, und die südliche Declination um \emph{sechs Minuten} grösser, als in der Ephemeride, die in der Monatl. Corresp. October 1810 abgedruckt ist, und diese Unterschiede nahmen immer mehr zu, wie folgende Übersicht zeigt.
\Needspace{15\baselineskip}
\begin{center}Correction der Ephemeride der Pallas wegen der Störungen:\end{center}
\begin{center}\fontsize{10}{12}\selectfont\setlength{\tabcolsep}{4pt}\renewcommand{\arraystretch}{1.15}\begin{tabular}{llrr|r|r}
 &  &  &  & \shortstack{gerade\\Aufsteigung} & Declination\\\hline
1810. & Oct. & 26 & $0^{\text{u}}$ & $-45'$ & $+6'$\\
 & Nov. & 7 & 12 & $-52$ & $+7$\\
 & --- & 20 & 0 & $-59$ & $+9$\\
 & Dec. & 2 & 12 & $-67$ & $+11$\\
 & --- & 15 & 0 & $-75$ & $+12$\\
 & --- & 27 & 12 & $-84$ & \ObservedDeclinationMark\,$13$\\
1811 & Jan. & 9 & 0 & $-94$ & $+14$\\
 & --- & 21 & 12 & $-103$ & $+14$\\
 & Febr. & 3 & 0 & $-112$ & $+14$\\
\end{tabular}\end{center}

\par\indent Ein so bedeutender Unterschied war für die Aufsuchung der Pallas von der grössten Wichtigkeit, und Hr. Prof. \textsc{Gauss} theilte deshalb diese Resultate sofort mehreren auswärtigen Astronomen mit, da das ungewöhnlich ungünstige Wetter hier die Aufsuchung nicht verstattete. Um indess in dieser wichtigen Angelegenheit ganz sicher zu gehen, hielt Hr. Prof. \textsc{Gauss} es für gut, die ganze Rechnung für die Perturbationen nochmals zu wiederholen, theils weil es so leicht möglich war, dass bei dieser weitläufigen Arbeit hier und da Fehler sich eingeschlichen haben konnten, theils weil schon an sich von der zweiten Rechnung, bei welcher schon die verbesserte Theorie zum Grunde gelegt werden konnte, noch zuverlässigere Resultate sich erwarten liessen. Inzwischen fanden sich nur unbedeutende Unterschiede von wenigen Secunden, die Uebereinstimmung der sechs Oppositionen blieb eben so gut, und die nothwendige Aenderung der Ephemeride bis auf Kleinigkeiten eben so gross. Um die jetzige Uebereinstimmung der sechs Oppositionen desto besser schätzen, und mit den grossen Unterschieden bei rein elliptischen Elementen (s. unsere Anz. 1810 S. 1971) [S. 63 d. B.] vergleichen zu können, setzen wir dieselbe hier her:
\Needspace{10\baselineskip}
\begin{center}\fontsize{10}{12}\selectfont\setlength{\tabcolsep}{4pt}\renewcommand{\arraystretch}{1.15}\begin{tabular}{r|r|r}
 & \multicolumn{2}{c}{Unterschied der}\\
 & \shortstack{mittlere\\Länge} & \shortstack{heliocentr.\\Breite}\\\hline
1803 & $+7''3$ & $+1''6$\\
1804 & $-6.8$ & $-4.0$\\
1805 & $+1.5$ & $-4.1$\\
1807 & $-3.3$ & $-3.8$\\
1808 & $-3.9$ & $-16.9$\\
1809 & $+5.0$ & $+3.5$\\
\end{tabular}\end{center}

\SourceLeaf{329}\par\indent Die Breite 1808 war bekanntlich schlecht beobachtet. Man ist berechtigt, hieraus zu schliessen, dass die Störungen, welche die andern Planeten, besonders Saturn und Mars, auf die Bewegung der Pallas ausüben, gegen die Störungen vom Jupiter ganz unbedeutend sind, und vor der Hand noch ganz vernachlässigt werden dürfen.

Auch die auswärtigen Astronomen, denen diese Resultate mitgetheilt waren, wurden von der Witterung nicht begünstigt; nur Hr. v. \textsc{Lindenau} beobachtete am 19. Dec. auf der Seeberger Sternwarte einen kleinen Stern am Passage-Instrumente.
\[\begin{gathered}16^{\text{u}}\;9^{\text{m}}\;7^{\text{s}}1\ \text{M. Z.}\quad\text{Ger. Aufst. }150^\circ\;20'\;12''45\\\text{geschätzte Declination }20^\circ\;58'\text{ südl.}\end{gathered}\]
{\rightskip=0pt plus 1em \noindent also auf die Minute übereinstimmend mit dem Platze, wo, nach der verbesserten Ephemeride, die Pallas stehen musste. Inzwischen gab diese einzelne Beobachtung, welche nachher noch nicht hatte bestätigt werden können, noch keine völlige Gewissheit, ob dies auch wirklich die Pallas gewesen sei.\par}

\par\indent Auf der hiesigen Sternwarte war die Nacht vom 4. zum 5. Januar nach einem Zwischenraume von beinahe zwei Monaten die erste, wo sich die Aufsuchung der Pallas versuchen liess. Da Hr. Prof. \textsc{Gauss} dazu nur bewegliche Instrumente anwenden konnte, so wurde das Aufsuchen durch den Mondschein und die strenge Kälte (11 Grad) ungemein erschwert. Indessen sah doch Hr. Prof. \textsc{Gauss} seine Anstrengungen von einem glücklichen Erfolge belohnt; er fand nicht allein die Pallas wieder auf dem Platze, wo er sie nach seiner verbesserten Ephemeride erwartete, sondern es gelangen ihm auch ein paar Vergleichungen mit Sternen der Histoire céleste, woraus er, mit Vorbehalt einer künftigen schärfern Reduction, folgende Stellung des Planeten ableitete:

\begin{center}\fontsize{10}{12}\selectfont\setlength{\tabcolsep}{4pt}\renewcommand{\arraystretch}{1.15}\begin{tabular}{c|c|c}
1811 Mittl. Zeit & gerade Aufsteigung & südl. Abweichung\\\hline
Jan. 4 $15^{\text{u}}\;4^{\text{m}}\;42^{\text{s}}$ & $151^\circ\;40'\;5''$ & $22^\circ\;22'\;36''$\\
\end{tabular}\end{center}

\par\indent Der Planet hatte übrigens erst die 8 bis 9. Grösse.

\Needspace{17\baselineskip}{\rightskip=0pt plus 1em \par\indent Auswärtige Astronomen werden vermittelst der eben angeführten Verbesserung der Ephemeride den Planeten leicht auffinden; will man dieselbe weiter ausdehnen, oder schärfer berechnen, so kann man sich dazu folgender Elemente bedienen, welche gegenwärtig die Bewegung des Planeten darstellen:\par}

\begin{center}\fontsize{10}{12}\selectfont\setlength{\tabcolsep}{4pt}\renewcommand{\arraystretch}{1.15}\begin{tabular}{llr}
\shortstack[l]{Epoche der mittlern Länge\\im Meridian von Göttingen 1811} & .\;. & $126^\circ\;33'\;6''$\\
Länge der Sonnennähe & .\;.\;.\;.\;.\;.\;.\;.\;.\;.\;.\;.\;.\;.\;.\;. & $120\;55\;8$\\
Länge des aufsteigenden Knoten & .\;.\;.\;.\;.\;.\;.\;.\;.\;.\;.\;. & $172\;32\;40$\\
Neigung der Bahn & .\;.\;.\;.\;.\;.\;.\;.\;.\;.\;.\;.\;.\;.\;.\;.\;. & $34\;35\;14$\\
Tägliche mittlere tropische Bewegung & .\;.\;.\;.\;.\;.\;.\;,\;.\;. & $768''980$\\
Excentricität & .\;.\;.\;.\;.\;.\;.\;.\;.\;.\;.\;.\;.\;.\;.\;.\;.\;.\;.\;. & $0.24160$\\
Logarithm der halben grossen Axe & .\;.\;.\;.\;.\;.\;.\;.\;.\;.\;.\;. & $0.44278$\\
\end{tabular}\end{center}

\par\indent Hr. Prof. \textsc{Gauss} glaubt durch die ungesäumte Bekanntmachung dieser Resultate den practischen Astronomen, die auf die Beobachtung der Pallas ausgehen wollen, einen Dienst zu erweisen, und behält sich eine vollständige Darlegung seiner Arbeit über die Störungen der Pallas auf eine andere Gelegenheit vor.
\clearpage\fancyhead[L]{\small GAUSS AN VON LINDENAU.}\SourceLeaf{330}
\begin{center}\textbf{GAUSS AN VON LINDENAU.}\par\vspace{4pt}\hrule\vspace{4pt}\fontsize{10}{12}\selectfont Monatliche Correspondenz. Band XXIII. Seite 97. 98. 1811. Januar.\par\vspace{4pt}\hrule\vspace{4pt}\end{center}

\begin{flushright}Göttingen, den 14. Januar 1811.\end{flushright}
\par\indent --- --- --- Auf der hiesigen Sternwarte war die Nacht vom 4. zum 5. Januar, nach einem Zwischenraum von beinahe zwei Monaten, die erste, wo sich die Aufsuchung der \emph{Pallas} versuchen liess. Da ich dazu nur bewegliche Instrumente anwenden konnte, so wurde das Aufsuchen durch den Mondschein und die strenge Kälte (11$^\circ$ Réaum.) ungemein erschwert. Indessen sah ich meine Bemühungen von einem glücklichen Erfolge belohnt, indem ich die \emph{Pallas} nicht allein auf dem Platze, wo sie nach meiner verbesserten Ephemeride sich befinden musste, auffand, sondern es gelangen mir auch ein paar Vergleichungen mit Sternen der \emph{Histoire céleste}, woraus ich mit Vorbehalt einer künftigen schärfern Reduction, folgende Stellung des Planeten ableitete:

\begin{center}\fontsize{10}{12}\selectfont\setlength{\tabcolsep}{4pt}\renewcommand{\arraystretch}{1.15}\begin{tabular}{c|c|c|c}
1811 & Mittlere Zeit & \shortstack{gerade Aufsteig.\\der Pallas} & \shortstack{Südliche\\Declination}\\\hline
4 Januar & $15^{\text{u}}\;4^{\text{m}}\;42^{\text{s}}$ & $151^\circ\;40'\;5''$ & $22^\circ\;22'\;36''$\\
\end{tabular}\end{center}

\Needspace{16\baselineskip}\par\indent Meine neu revidirten Störungsrechnungen haben mir nun folgende Elemente der \emph{Pallas} gegeben, die von den Ihnen zuletzt mitgetheilten (Dec. Heft 1810 S. 592) [S. 321 d. B.] nur unbedeutend abweichen.

\begin{center}\fontsize{10}{12}\selectfont\setlength{\tabcolsep}{4pt}\renewcommand{\arraystretch}{1.15}\begin{tabular}{llr}
\shortstack[l]{Epoche der mittl. Länge\\im Meridian von Göttingen} & .\;.\;. & $126^\circ\;33'\;6''$\\
Länge der Sonnennähe & .\;.\;.\;.\;.\;.\;.\;.\;.\;.\;.\;.\;.\;.\;. & $120\;55\;8$\\
Länge des aufsteigenden Knotens & .\;.\;.\;.\;.\;.\;.\;.\;.\;.\;. & $172\;32\;40$\\
Neigung der Bahn & .\;.\;.\;.\;.\;.\;.\;.\;.\;.\;.\;.\;.\;.\;.\;. & $34\;35\;14$\\
Tägliche mittlere tropische Bewegung & .\;.\;.\;.\;.\;.\;.\;.\;. & $768.980$\\
Excentricität & .\;.\;.\;.\;.\;.\;.\;.\;.\;.\;.\;.\;.\;.\;.\;.\;.\;. & $0.24160$\\
Logarithmus der halben grossen Axe & .\;.\;.\;.\;.\;.\;.\;.\;. & $0.44278$\\
\end{tabular}\end{center}

\par\indent Hier ist noch die berechnete Fortsetzung der Ephemeride: [vom 3. Febr. bis 12. April 1811].
